% Appendix D -- Computation (draft prose; from the R0 code and records)
%   code      model-codes/codes/epc-lifecycle-deep-learning/notes/
%             model-assumptions-review-2026-09-22/parallel-diagnostics-2026-09-25/
%             completed-backup/code-training-first-full-20260925 (frozen R0 code)
%             epc_backward/training_first/workflow.py (TrainingConfig, train_stage),
%             epc_backward/nets.py (ValueNet baseline), physics.py (controls),
%             features.py (inputs), simulate.py
%   sizes     value net 42-512-512-512-1 (checkpoint tensors), policy net
%             42-128-128-57; 20,480 rows per age; 71 stages, 9.71 h
%   accuracy  replication/data/model/r0/final-validation/final-estimate-euler.json
%             (128 states per age; euler_consumption_fraction_abs, interior states)
%   search    replication/data/model/r0/estimation/estimation_results.json
%   SEs       replication/code/3-model/standard_errors.py;
%             replication/data/model/standard-errors-n4096/summary.json
\section{Computation}\label{app:computation}

\subsection{Networks.}\label{app:computation-networks}
Each age has its own value network and policy network, 120 networks for the
60 ages. Both take the same 42 inputs: the state, functions of it, the
seven parameters in $\vartheta$, and the discount factor $\beta$. Each parameter is scaled to $[0,1]$ over its
search range, the bequest scale $\upsilon$ on a log scale.

The value network has three hidden layers of 512 units with the SiLU
activation. The activation is smooth because the maximization at the
previous age reads the slope of the value in liquid wealth. The network does
not fit the value directly. When the bequest motive is strong, the value of
a household is dominated by the bequest term, and the accuracy the choices
need is a small fraction of it. The network therefore fits the deviation
from an analytic baseline,
\begin{equation*}
  \widehat V_j(\mathbf x_j,\vartheta)=A_j(\beta)\,g(\tilde q_j)
  +\bar v_j+\varsigma_j\,N_j(\mathbf x_j,\vartheta),
  \qquad
  A_j(\beta)=\sum_{k\geq0}\beta^k\Bigl(\prod_{i=0}^{k-1}\phi_{j+i}\Bigr)
  (1-\phi_{j+k}),
\end{equation*}
where $g$ is the bequest utility~\eqref{eq:bequest}, $\tilde q_j$ is the
household's net worth at the start of the period (cash on hand plus home
equity), $A_j(\beta)$ is the discounted probability of dying at each future
age, $N_j$ is the network, and $\bar v_j$ and $\varsigma_j$ standardize the values at
age $j$.

The policy network has two hidden layers of 128 units and two to four
outputs for each alternative. For every alternative it predicts the
logarithm of consumption and the probability that the household consumes
the most the alternative allows: it saves nothing and, if it buys, borrows
only what the purchase requires. When this probability exceeds 0.9, the
household consumes that maximum. An owner who stays also has a payment
share, which places the mortgage payment between the smallest payment the
contract allows and the largest compatible with the predicted consumption.
A buyer also has the size share $a_h$ and the mortgage share $a_m$,
described below. A renter's home size is the optimal size given total
spending, in closed form, because it does not affect the future. Liquid
saving is the residual in every alternative.

These outputs are not the choices themselves. A network outputs
unrestricted real numbers, but the choices must satisfy the budget and
borrowing constraints, so a fixed mapping turns the outputs into a feasible
allocation. Take a renter who buys a home. The network predicts
consumption $c$ and the two shares $a_h,a_m\in[0,1]$. The first places the
size of the new home in the feasible range,
\begin{equation*}
  h'=\underline h^{\mathrm{own}}
  +a_h\bigl(\overline h^{\mathrm{own}}-\underline h^{\mathrm{own}}\bigr),
\end{equation*}
and the second places the mortgage between the smallest balance that
finances the purchase together with $c$ and the largest balance that the
loan-to-value and payment limits allow. Liquid saving is what is left, so
the budget holds by construction. If $c$ is not affordable at $h'$, the
size is reduced until it is. An error in $a_m$ then moves money between the
mortgage and liquid saving and leaves consumption unchanged. A purchase
involves large sums relative to a year's consumption, so this matters: if
the network predicted saving instead, a small error in it would be a large
error in consumption.

\subsection{Training.}\label{app:computation-training}
At each age, the training sample has 20,480 points. It consists of 4,096
draws of the state and the parameters, each with four values of liquid
wealth, and 512 groups of points with small mortgages. The parameters are
drawn uniformly over the box, $\upsilon$ log-uniformly. Ten percent of the
draws are held out to measure the fit.

At each point the expectation in equation~\eqref{eq:end-value} is computed
by quadrature, with five nodes for the permanent income shock and three for
the transitory shock, and exactly over whether an efficient home stays
efficient. For every feasible alternative, the maximization in
equation~\eqref{eq:choice-value} combines a grid in the choices refined by
golden-section search, a grid in consumption, and candidates on the
boundary of the feasible set, and keeps the best. The target of the value
network is the logsum~\eqref{eq:logsum}. By the envelope condition, its
derivative with respect to liquid wealth is $(1+r_b)$ times the
probability-weighted marginal utility of consumption across alternatives.

The value network minimizes the relative error in the value plus 0.01 times
the error in this derivative, in 24,000 steps of the Adam optimizer, with a
learning rate that decays from $10^{-4}$ to $10^{-6}$. The policy network
fits the maximizing choices in 8,000 steps, with a learning rate that
decays from $10^{-3}$ to $10^{-5}$. In its loss, the choices of each
alternative are weighted by $0.05+0.95$ times the probability that the
household chooses it, so the fit concentrates on the alternatives that
households choose. Infeasible alternatives have zero weight. Both networks
start from the trained networks of age $j+1$.

The networks have no explicit penalty on their weights, such as weight
decay or dropout. Instead, three devices keep them from fitting noise in
the sample. Every 2,000 steps for the value network and every 1,000 steps
for the policy network, the loss is computed on the held-out points, and
the weights with the lowest held-out loss are kept. The norm of the
gradient is clipped at 5. And each step uses a random minibatch of 2,048
points. Table~\ref{tab:networks} collects these settings.

\begin{table}[htbp]
  \centering
  \caption{Network architecture and training}\label{tab:networks}
  \begin{threeparttable}
    \small
    \begin{tabular}{@{}l>{\raggedright\arraybackslash}p{0.33\textwidth}%
                       >{\raggedright\arraybackslash}p{0.33\textwidth}@{}}
      \toprule
      & Value network & Policy network \\
      \midrule
      \multicolumn{3}{@{}l}{\textit{Architecture}} \\
      \quad Inputs & 42 & 42 \\
      \quad Hidden layers & $3\times512$ units & $2\times128$ units \\
      \quad Activation & SiLU & SiLU \\
      \quad Outputs & 1, the deviation from the bequest baseline
        & 57, two to four for each alternative \\
      \addlinespace
      \multicolumn{3}{@{}l}{\textit{Sample at each age}} \\
      \quad Points & 20,480 & 20,480 \\
      \quad Held out & 10 percent of the draws & 10 percent of the draws \\
      \addlinespace
      \multicolumn{3}{@{}l}{\textit{Loss}} \\
      \quad Fit & squared error relative to $1+|V|$
        & Huber ($\delta=0.25$) in log consumption; cross-entropy in shares
          and probabilities \\
      \quad Derivative & $0.01\times$ squared relative error in
        $\partial V/\partial b$ & -- \\
      \quad Weight of a point & equal
        & $0.05+0.95\times$ choice probability of the alternative \\
      \addlinespace
      \multicolumn{3}{@{}l}{\textit{Optimization}} \\
      \quad Optimizer & Adam & Adam \\
      \quad Steps & 24,000 & 8,000 \\
      \quad Learning rate, cosine decay & $10^{-4}$ to $10^{-6}$
        & $10^{-3}$ to $10^{-5}$ \\
      \quad Minibatch & 2,048 points & 2,048 points \\
      \quad Initial weights & networks of age $j+1$ & networks of age $j+1$ \\
      \addlinespace
      \multicolumn{3}{@{}l}{\textit{Regularization}} \\
      \quad Weight decay, dropout & none & none \\
      \quad Gradient-norm clipping & 5 & 5 \\
      \quad Early stopping & best held-out loss, checked every 2,000 steps
        & best held-out loss, checked every 1,000 steps \\
      \bottomrule
    \end{tabular}
    \begin{tablenotes}[flushleft]\footnotesize
      \item \textit{Notes:} Each of the 60 ages has one network of each
      kind. The sample has 4,096 draws of the state and the parameters, each
      at four values of liquid wealth, and 512 draws of owners, each with a
      small mortgage and with none, at the same four values. Points that
      share a draw are held out together.
    \end{tablenotes}
  \end{threeparttable}
\end{table}

At ages 82 and above the problem is stationary. The networks of this age
are fitted to targets computed from their previous iteration, at least 6
and at most 12 times. The iteration stops early when the choices change by
less than a tolerance in two consecutive iterations.
With the 59 younger ages, training has 71 stages. They took 9.7
hours on one GPU, 8.8 of them in computing the targets.

\subsection{Accuracy.}\label{app:computation-accuracy}
On the held-out points at age 23, the value network's mean relative error
is 0.3 percent and its 95th percentile is 1.0 percent. We also check the
policies at the estimate with Euler-equation errors, expressed as a share of
consumption, at 128 states at each of the ages 23, 40, 64, 65, 78, and 82+,
among the states in which neither the borrowing limit nor the consumption
floor binds. The median error is 0.9 percent at 23 and at 82+, 1.2 percent at
78, 3.3 percent at 40, and 2.9 and 4.7 percent at 64 and 65. The 95th
percentiles range from 3.6 percent at 78 to 13.6 percent at 65. The largest
error, 54.6 percent, is at age 40. The errors are largest around
retirement, where income falls.

\subsection{Search.}\label{app:computation-search}
The search has three stages, all with the same draws of the shocks. First,
it simulates 96 points of a scrambled Sobol sequence over the box with
2,048 households each. Three of the points are replaced by estimates from
earlier runs. It keeps the best 12. Second, it refines them locally over
three rounds, drawing 16, 12, and 8 normal perturbations of each candidate
with standard deviations of 10, 5, and 2 percent of the width of each
range, with 4,096 households, and keeping the best 8 after each round.
Third, it simulates the best 10 candidates with 32,768 households, and the
best of them is the estimate. The search evaluates 486 candidates. The
moments in Panel~B of Table~\ref{tab:estimates} are simulated at the estimate with 65,536
households and new draws.


\subsection{Standard errors.}\label{app:computation-se}
The derivative $G$ in equation~\eqref{eq:smm-variance} is computed by
central differences at $\pm3$ percent of the width of each range, with
65,536 households and the same draws at every point, on the scale on which
the moments enter the criterion: levels for shares, logs for ratios. The
covariance $\Omega$ of the data moments comes from a bootstrap of the
surveys that resamples clusters, on the same scale.

The simulation noise is $\Sigma_S=r\,D\Omega D$, with
$D=\operatorname{diag}(1/\sqrt{S_k})$. Here $S_k$ is the number of simulated
households in the age group of moment $k$, at the 4,096 households of the
search, divided by the number of observations in the data. A simulated
household counts once however many ages it contributes, so $D\Omega D$
overstates the noise. The factor $r=0.069$ corrects for this. It is the
ratio of the simulation noise measured at four parameter vectors, as the
difference between two independent sets of draws with 65,536 households, to
the noise $D\Omega D$ predicts for that number of households.
